\par\Needspace{6\baselineskip}\medskip
\hypertarget{native-121-article090}{}\pdfbookmark[2]{90.}{article-090}
\begin{center}90.\end{center}
Considerabimus primo casum eum, ubi e solutione aequationis 12 valor non nimis magnus ipsius $g$ emergit, ita ut $\dfrac{2g-\sin 2g}{\sin g^3}$ in seriem secundum potestates ipsius $\sin\tfrac12 g$ progredientem evolvere liceat. Numerator huius expressionis, quam per $X$ denotabimus, fit
\[
=\dfrac{32}{3}\sin\tfrac12 g^3-\dfrac{16}{5}\sin\tfrac12 g^5-\dfrac47\sin\tfrac12 g^7-\text{etc.}
\]
Denominator autem
\[
=8\sin\tfrac12 g^3-12\sin\tfrac12 g^5+3\sin\tfrac12 g^7+\text{etc.}
\]
Unde $X$ obtinet formam
\[
\dfrac43+\dfrac85\sin\tfrac12 g^2+\dfrac{64}{35}\sin\tfrac12 g^4+\text{etc.}
\]
Ut autem legem progressionis coëfficientium eruamus, differentiamus aequationem $X\sin g^3=2g-\sin 2g$, unde prodit
\[
3X\cos g\sin g^2+\sin g^3\dfrac{\mathrm{d}X}{\mathrm{d}g}=2-2\cos 2g=4\sin g^2;
\]
statuendo porro $\sin\tfrac12 g^2=x$, fit $\dfrac{\mathrm{d}x}{\mathrm{d}g}=\tfrac12\sin g$, unde concluditur
\[
\dfrac{\mathrm{d}X}{\mathrm{d}x}=\dfrac{8-6X\cos g}{\sin g^2}=\dfrac{4-3X(1-2x)}{2x(1-x)}\mathpunct{\raisebox{.45ex}{,}}
\]
et proin
\[
(2x-2xx)\dfrac{\mathrm{d}X}{\mathrm{d}x}=4-(3-6x)X.
\]
\SourceLeaf{122}Quodsi igitur statuimus
\[
X=\tfrac43(1+\alpha x+\beta xx+\gamma x^3+\delta x^4+\text{etc.}),
\]
obtinemus aequationem
\[
\begin{gathered}
\tfrac83\{\alpha x+(2\beta-\alpha)xx+(3\gamma-2\beta)x^3+(4\delta-3\gamma)x^4+\text{etc.}\}
\\[4pt]
=(8-4\alpha)x+(8\alpha-4\beta)xx+(8\beta-4\gamma)x^3+(8\gamma-4\delta)x^4+\text{etc.},
\end{gathered}
\]
quae identica esse debet. Hinc colligimus $\alpha=\tfrac65$, $\beta=\tfrac87\alpha$, $\gamma=\tfrac{10}{9}\beta$, $\delta=\tfrac{12}{11}\gamma$ etc., ubi lex progressionis obvia est. Habemus itaque
\[
\begin{aligned}
X={}&\tfrac43+\dfrac{4\mathbin{.}6}{3\mathbin{.}5}x+\dfrac{4\mathbin{.}6\mathbin{.}8}{3\mathbin{.}5\mathbin{.}7}xx+\dfrac{4\mathbin{.}6\mathbin{.}8\mathbin{.}10}{3\mathbin{.}5\mathbin{.}7\mathbin{.}9}x^3
\\[4pt]
&+\dfrac{4\mathbin{.}6\mathbin{.}8\mathbin{.}10\mathbin{.}12}{3\mathbin{.}5\mathbin{.}7\mathbin{.}9\mathbin{.}11}x^4+\text{etc.}
\end{aligned}
\]

Hanc seriem transformare licet in fractionem continuam sequentem:
\[
X=\cfrac{\tfrac43}{1-\cfrac{\tfrac65x}{1+\cfrac{\tfrac{2}{5\mathbin{.}7}x}{1-\cfrac{\tfrac{5\mathbin{.}8}{7\mathbin{.}9}x}{1-\cfrac{\tfrac{1\mathbin{.}4}{9\mathbin{.}11}x}{1-\cfrac{\tfrac{7\mathbin{.}10}{11\mathbin{.}13}x}{1-\cfrac{\tfrac{3\mathbin{.}6}{13\mathbin{.}15}x}{1-\cfrac{\tfrac{9\mathbin{.}12}{15\mathbin{.}17}x}{1-\text{etc.}}}}}}}}}
\]
Lex secundum quam coëfficientes $\tfrac65$, $-\tfrac{2}{5\mathbin{.}7}$, $\tfrac{5\mathbin{.}8}{7\mathbin{.}9}$, $\tfrac{1\mathbin{.}4}{9\mathbin{.}11}$ etc. progrediuntur obvia est; scilicet terminus $n^{\mathrm{tus}}$ huius seriei fit pro $n$ pari $=\dfrac{(n-3)n}{(2n+1)(2n+3)}$, pro $n$ impari autem $=\dfrac{(n+2)(n+5)}{(2n+1)(2n+3)}$: ulterior huius argumenti evolutio\HistoricalNoteEight{Vgl. Band III. S. 135.]} nimis aliena esset ab instituto nostro. Quodsi iam statuimus
\[
\cfrac{x}{1+\cfrac{\tfrac{2}{5\mathbin{.}7}x}{1-\cfrac{\tfrac{5\mathbin{.}8}{7\mathbin{.}9}x}{1-\cfrac{\tfrac{1\mathbin{.}4}{9\mathbin{.}11}x}{1-\text{etc.}}}}}=x-\xi,
\]
\SourceLeaf{123}fit $X=\dfrac{1}{\frac34-\frac{9}{10}(x-\xi)}\mathpunct{\raisebox{.45ex}{,}}$ atque $\xi=x-\frac56+\frac{10}{9X}\mathpunct{\raisebox{.45ex}{,}}$ sive
\[
\xi=\dfrac{\sin g^3-\frac34(2g-\sin 2g)(1-\frac65\sin\frac12 g^2)}{\frac{9}{10}(2g-\sin 2g)}\mathpunct{\cdot}
\]

Numerator huius expressionis est quantitas ordinis septimi, denominator ordinis tertii, adeoque $\xi$ ordinis quarti, siquidem $g$ tamquam quantitas ordinis primi, sive $x$ tamquam ordinis secundi spectatur. Hinc concluditur, formulam hancce ad computum numericum exactum ipsius $\xi$ haud idoneam esse, quoties $g$ angulum non valde considerabilem exprimat: tunc autem ad hunc finem commode adhibentur formulae sequentes, quae ab invicem per ordinem commutatum numeratorum in coëfficientibus fractis differunt, et quarum prior e valore supposito ipsius $x-\xi$ haud difficile derivatur\footnote{Deductio posterioris quasdam transformationes minus obvias aliaque occasione [Band III. S. 137] explicandas supponit.
\par [Handschriftliche Bemerkung:] Fit
\[
\xi=\dfrac{2}{35}xx\dfrac{1+\frac{2\mathbin{.}8}{9}x+\frac{3\mathbin{.}8\mathbin{.}10}{9\mathbin{.}11}xx+\frac{4\mathbin{.}8\mathbin{.}10\mathbin{.}12}{9\mathbin{.}11\mathbin{.}13}x^3+\text{etc.}}{1+\frac65x+\frac{6\mathbin{.}8}{5\mathbin{.}7}xx+\frac{6\mathbin{.}8\mathbin{.}10}{5\mathbin{.}7\mathbin{.}9}x^3+\text{etc.}}\mathpunct{\cdot}
\]
Statuendo itaque $1+\frac{2\mathbin{.}8}{9}x+\frac{3\mathbin{.}8\mathbin{.}10}{9\mathbin{.}11}xx+\text{etc.}=A$, fit $\xi=\dfrac{\frac{2}{35}Axx(1-\frac65x)}{1-\frac{12}{175}Axx}\mathpunct{\cdot}$ Habetur etiam
\[
A=(1-x)^{-\frac32}\left(1+\frac{1\mathbin{.}5}{2\mathbin{.}9}x+\frac{1\mathbin{.}3\mathbin{.}5\mathbin{.}7}{2\mathbin{.}4\mathbin{.}9\mathbin{.}11}xx+\text{etc.}\right).
\]
}:
\[
\mathrm{13)}\qquad \xi=\cfrac{\tfrac{2}{35}xx}{1+\tfrac{2}{35}x-\cfrac{\tfrac{40}{63}x}{1-\cfrac{\tfrac{4}{99}x}{1-\cfrac{\tfrac{70}{143}x}{1-\cfrac{\tfrac{18}{195}x}{1-\cfrac{\tfrac{108}{255}x}{1-\text{etc.}}}}}}}
\]
sive
\[
\xi=\cfrac{\tfrac{2}{35}xx}{1-\tfrac{18}{35}x-\cfrac{\tfrac{4}{63}x}{1-\cfrac{\tfrac{40}{99}x}{1-\cfrac{\tfrac{18}{143}x}{1-\cfrac{\tfrac{70}{195}x}{1-\cfrac{\tfrac{40}{255}x}{1-\text{etc.}}}}}}}
\]
\SourceLeaf{124}In tabula tertia huic operi annexa pro cunctis valoribus ipsius $x$ a 0 usque ad 0,3, per singulas partes millesimas, valores respondentes ipsius $\xi$ ad septem figuras decimales computati reperiuntur. Haec tabula primo aspectu monstrat exiguitatem ipsius $\xi$ pro valoribus modicis ipsius $g$; ita e. g. pro $E'-E=10^0$, sive $g=5^0$, ubi $x=0{,}00190$, fit $\xi=0{,}000\ 0002$. Superfluum fuisset, tabulam adhuc ulterius continuare, quum termino ultimo $x=0{,}3$ respondeat $g=66^0 25'$ sive $E'-E=132^0 50'$. Ceterum tabulae columna tertia, quae valores ipsius $\xi$ valoribus negativis ipsius $x$ respondentes continet, infra loco suo explicabitur.
